\chapter{Transfer Matrix and Quantum--Classical Correspondence}
\label{chap:transfer-matrix-quantum-classical}

The previous chapters developed the free-fermion machinery used to diagonalize
one-dimensional quantum spin chains: Jordan--Wigner transformation, Fourier
transform, BdG blocks, Bogoliubov quasiparticles, and Gaussian correlation
functions. Transfer matrices provide a complementary route to exact solvability.
They are indispensable for two related reasons:
\begin{enumerate}[label=(\roman*)]
    \item a classical statistical model in one higher dimension can often be
    interpreted as an imaginary-time path integral of a quantum model; and
    \item a family of commuting transfer matrices can generate a quantum
    Hamiltonian by a logarithmic derivative or Hamiltonian limit.
\end{enumerate}
The first point is the quantum--classical correspondence. The second point is the
bridge to vertex-model and Bethe-ansatz integrability in later chapters.

\section{Classical transfer matrices}
\label{sec:classical-transfer-matrices}

Consider a classical statistical model on a strip with $M$ layers. Let
$\bm{s}_m$ denote the entire configuration on the $m$th layer. If the Boltzmann
weight couples only neighbouring layers, the partition function can be written as
\begin{equation}
\label{eq:transfer-matrix-general-partition-function}
    Z
    =
    \sum_{\bm{s}_1,\ldots,\bm{s}_M}
    T_{\bm{s}_1\bm{s}_2}
    T_{\bm{s}_2\bm{s}_3}
    \cdots
    T_{\bm{s}_M\bm{s}_1}
    =
    \Tr T^M .
\end{equation}
The matrix $T$ is the \emph{transfer matrix}. Its matrix element
$T_{\bm{s}\bm{s}'}$ is the Boltzmann weight for two adjacent layers
$\bm{s}$ and $\bm{s}'$, including a chosen assignment of intra-layer weights.
The trace implements periodic boundary conditions in the transfer direction.

If $T$ is finite-dimensional and diagonalizable, with eigenvalues
\begin{equation}
    \lambda_0,\lambda_1,\lambda_2,\ldots,
    \qquad
    \abs{\lambda_0}\geq \abs{\lambda_1}\geq\cdots,
\end{equation}
then
\begin{equation}
\label{eq:transfer-matrix-free-energy}
    Z=\sum_n \lambda_n^M .
\end{equation}
In the thermodynamic limit $M\to\infty$, the largest eigenvalue controls the free
energy per transfer step. If each layer contains $L$ sites, then
\begin{equation}
\label{eq:transfer-matrix-free-energy-density}
    f
    =
    -\frac{1}{\beta_{\rm cl}L}
    \lim_{M\to\infty}\frac{1}{M}\log Z
    =
    -\frac{1}{\beta_{\rm cl}L}\log \lambda_0,
\end{equation}
up to the convention-dependent normalization of the classical inverse temperature
$\beta_{\rm cl}$.

Transfer matrices also encode correlation lengths. For a layer observable $O$,
\begin{equation}
\label{eq:transfer-matrix-correlation-spectral}
    \langle O_m O_0\rangle
    =
    \frac{1}{Z}\Tr\left(O T^m O T^{M-m}\right).
\end{equation}
After inserting the spectral resolution of $T$, the connected correlation
function has the large-distance form
\begin{equation}
\label{eq:transfer-matrix-correlation-length}
    \langle O_m O_0\rangle_c
    \sim
    \left(\frac{\lambda_1}{\lambda_0}\right)^m,
    \qquad
    \xi^{-1}
    =
    \log\frac{\lambda_0}{\abs{\lambda_1}},
\end{equation}
provided $O$ has nonzero matrix element between the leading eigenstate and the
first relevant subleading eigenstate. Thus the spectrum of $T$ plays the same
role for a classical model that the spectrum of a Hamiltonian plays for a quantum
model.

\section{Warm-up: the one-dimensional classical Ising chain}
\label{sec:transfer-matrix-1d-classical-ising}

For the one-dimensional classical Ising model with spins $s_j=\pm1$ and periodic
boundary conditions,
\begin{equation}
\label{eq:classical-ising-1d-partition-function}
    Z_L(K)
    =
    \sum_{s_1,\ldots,s_L=\pm1}
    \exp\left(K\sum_{j=1}^{L}s_js_{j+1}\right),
    \qquad
    s_{L+1}=s_1,
\end{equation}
the transfer matrix is the $2\times2$ matrix
\begin{equation}
\label{eq:1d-ising-transfer-matrix}
    T_{ss'}=\ee^{Kss'},
    \qquad
    T=
    \begin{pmatrix}
        \ee^K & \ee^{-K} \\
        \ee^{-K} & \ee^K
    \end{pmatrix}.
\end{equation}
Its eigenvalues are
\begin{equation}
\label{eq:1d-ising-transfer-eigenvalues}
    \lambda_+=\ee^K+\ee^{-K}=2\cosh K,
    \qquad
    \lambda_- =\ee^K-\ee^{-K}=2\sinh K.
\end{equation}
Therefore
\begin{equation}
\label{eq:1d-ising-exact-partition-function}
    Z_L(K)=\lambda_+^L+\lambda_-^L.
\end{equation}
For finite $K$, the correlation length is finite:
\begin{equation}
\label{eq:1d-ising-correlation-length}
    \xi^{-1}
    =
    \log\frac{\lambda_+}{\lambda_-}
    =
    \log\coth K .
\end{equation}
This elementary example already contains the main lesson: thermodynamics and
correlations are spectral properties of the transfer matrix.

\section{Transfer matrix as imaginary-time evolution}
\label{sec:transfer-matrix-imaginary-time-evolution}

A quantum thermal partition function is
\begin{equation}
\label{eq:quantum-thermal-partition-function}
    Z_{\rm q}(\beta)
    =
    \Tr \ee^{-\beta H} .
\end{equation}
If imaginary time is discretized into $M$ steps,
\begin{equation}
    \beta=M\delta\tau,
\end{equation}
then
\begin{equation}
\label{eq:quantum-partition-transfer-matrix}
    Z_{\rm q}(\beta)
    =
    \Tr\left(\ee^{-\delta\tau H}\right)^M .
\end{equation}
Thus the short imaginary-time evolution operator
\begin{equation}
\label{eq:quantum-transfer-matrix-definition}
    T_{\rm q}(\delta\tau)=\ee^{-\delta\tau H}
\end{equation}
is itself a transfer matrix. Conversely, whenever a positive transfer matrix $T$
can be written as
\begin{equation}
\label{eq:hamiltonian-from-transfer-matrix}
    T=C\ee^{-a_\tau H},
\end{equation}
with $C>0$ and $a_\tau$ interpreted as an imaginary-time lattice spacing, it
defines the quantum Hamiltonian
\begin{equation}
\label{eq:hamiltonian-log-transfer-matrix}
    H
    =
    -\frac{1}{a_\tau}\log T
    +
    \frac{\log C}{a_\tau}\id .
\end{equation}
The additive constant proportional to $\id$ shifts the zero of energy and does not
affect normalized expectation values or energy gaps.

If $T\ket{n}=\lambda_n\ket{n}$ and
$T=C\ee^{-a_\tau H}$, then
\begin{equation}
\label{eq:transfer-eigenvalues-energy-levels}
    \lambda_n=C\ee^{-a_\tau E_n},
    \qquad
    E_n-E_0
    =
    -\frac{1}{a_\tau}
    \log\frac{\lambda_n}{\lambda_0} .
\end{equation}
Therefore transfer-matrix gaps become quantum energy gaps. In particular,
\begin{equation}
\label{eq:correlation-length-quantum-gap}
    \xi_\tau^{-1}
    =
    \log\frac{\lambda_0}{\abs{\lambda_1}}
    =
    a_\tau\,(E_1-E_0)
\end{equation}
when the first subleading eigenvalue controls the decay.

\begin{remark}[Hermiticity and positivity]
A symmetric positive transfer matrix defines a Hermitian Hamiltonian by
$H=-a_\tau^{-1}\log T$, up to a constant. More general transfer matrices may be
non-Hermitian. Such matrices are still useful, for example in stochastic and
non-equilibrium problems, but the direct interpretation as a Hermitian quantum
Hamiltonian then requires additional structure.
\end{remark}

\section{Trotter decomposition and the extra dimension}
\label{sec:trotter-extra-dimension}

The relation between a $d$-dimensional quantum model and a $(d+1)$-dimensional
classical model follows from the Trotter product formula. Suppose
\begin{equation}
\label{eq:hamiltonian-split-A-B}
    H=A+B,
\end{equation}
where $A$ and $B$ do not commute but have simple matrix elements in appropriate
bases. Then
\begin{equation}
\label{eq:trotter-first-order}
    \ee^{-\beta H}
    =
    \lim_{M\to\infty}
    \left(
        \ee^{-\delta\tau A}\ee^{-\delta\tau B}
    \right)^M,
    \qquad
    \delta\tau=\frac{\beta}{M}.
\end{equation}
A symmetric version is
\begin{equation}
\label{eq:trotter-symmetric}
    \ee^{-\delta\tau(A+B)}
    =
    \ee^{-\frac{\delta\tau}{2}A}
    \ee^{-\delta\tau B}
    \ee^{-\frac{\delta\tau}{2}A}
    +O(\delta\tau^3).
\end{equation}
Inserting complete sets of states between the $M$ imaginary-time factors produces
a classical sum over configurations on $M$ time slices. The additional discrete
index is imaginary time. Schematically,
\begin{equation}
\label{eq:quantum-classical-schematic}
    Z_{d\text{-dimensional quantum}}
    =
    \Tr\ee^{-\beta H}
    \quad\longleftrightarrow\quad
    Z_{(d+1)\text{-dimensional classical}} .
\end{equation}
This correspondence is exact in the Trotter limit $M\to\infty$ and becomes a
controlled approximation at finite $M$.

The correspondence is not a statement that the quantum model and the classical
model are identical microscopic objects. Rather, their partition functions and
correlation functions are two representations of the same Euclidean statistical
sum. Quantum noncommutativity is converted into classical coupling along the
imaginary-time direction.

\section{From the anisotropic two-dimensional Ising model to the TFIM}
\label{sec:anisotropic-ising-to-tfim-transfer}

The most important example for these notes is the correspondence between the
anisotropic two-dimensional classical Ising model and the one-dimensional TFIM.
Let
\begin{equation}
\label{eq:anisotropic-2d-ising-partition-function}
    Z_{\rm 2D}
    =
    \sum_{s_{j,m}=\pm1}
    \exp\left[
        K_x\sum_{j,m}s_{j,m}s_{j+1,m}
        +
        K_\tau\sum_{j,m}s_{j,m}s_{j,m+1}
    \right],
\end{equation}
where $j=1,\ldots,L$ is the spatial coordinate and $m=1,\ldots,M$ is the transfer
direction. The couplings $K_x$ and $K_\tau$ are classical dimensionless
couplings.

To match the Pauli convention used in these notes, take the classical variable
$s_j=\pm1$ to label the eigenvalue of $\sigma_j^x$:
\begin{equation}
\label{eq:x-basis-transfer}
    \sigma_j^x\ket{s_j}_x=s_j\ket{s_j}_x .
\end{equation}
The horizontal classical coupling is diagonal in this basis and gives the
operator
\begin{equation}
\label{eq:horizontal-transfer-operator}
    \exp\left(K_x\sum_{j=1}^{L}\sigma_j^x\sigma_{j+1}^x\right).
\end{equation}
The vertical one-site Boltzmann weight is
\begin{equation}
\label{eq:vertical-weight-matrix}
    W_{ss'}=\ee^{K_\tau ss'}
    =
    \begin{pmatrix}
        \ee^{K_\tau} & \ee^{-K_\tau} \\
        \ee^{-K_\tau} & \ee^{K_\tau}
    \end{pmatrix}_{s,s'} .
\end{equation}
In the $\sigma^x$ eigenbasis, the operator $\sigma^z$ flips $s$. Hence
\begin{equation}
\label{eq:vertical-weight-exp-sigmaz}
    W
    =
    A\ee^{\kappa\sigma^z},
    \qquad
    \tanh\kappa=\ee^{-2K_\tau},
    \qquad
    A=\sqrt{2\sinh(2K_\tau)} .
\end{equation}
Indeed,
\begin{equation}
    A\cosh\kappa=\ee^{K_\tau},
    \qquad
    A\sinh\kappa=\ee^{-K_\tau}.
\end{equation}
A convenient symmetric row-to-row transfer matrix is therefore
\begin{equation}
\label{eq:2d-ising-row-transfer-operator}
    T
    =
    A^L
    \exp\left(
        \frac{K_x}{2}\sum_j\sigma_j^x\sigma_{j+1}^x
    \right)
    \exp\left(
        \kappa\sum_j\sigma_j^z
    \right)
    \exp\left(
        \frac{K_x}{2}\sum_j\sigma_j^x\sigma_{j+1}^x
    \right).
\end{equation}
Using the symmetric Trotter formula, the strongly anisotropic limit is obtained
by setting
\begin{equation}
\label{eq:ising-tfim-parameter-scaling}
    K_x=\delta\tau J,
    \qquad
    \kappa=\delta\tau h,
    \qquad
    \tanh(\delta\tau h)=\ee^{-2K_\tau}.
\end{equation}
Then
\begin{equation}
\label{eq:2d-ising-transfer-tfim-limit}
    T
    =
    A^L
    \exp\left[-\delta\tau H_{\rm TFIM}+O(\delta\tau^3)\right],
\end{equation}
with
\begin{equation}
\label{eq:transfer-chapter-tfim-hamiltonian}
    H_{\rm TFIM}
    =
    -J\sum_j\sigma_j^x\sigma_{j+1}^x
    -h\sum_j\sigma_j^z .
\end{equation}
This is precisely the TFIM convention fixed earlier in these notes.

For small $\delta\tau$,
\begin{equation}
\label{eq:large-vertical-coupling}
    \ee^{-2K_\tau}
    =
    \tanh(\delta\tau h)
    =
    \delta\tau h+O(\delta\tau^3),
\end{equation}
so the imaginary-time coupling $K_\tau$ is large. The quantum limit is therefore
a strongly anisotropic limit of the classical two-dimensional model:
\begin{equation}
    K_x\to0,
    \qquad
    K_\tau\to\infty,
    \qquad
    K_x\ee^{2K_\tau}\;\text{fixed}.
\end{equation}
The classical critical condition of the anisotropic two-dimensional Ising model,
\begin{equation}
\label{eq:2d-ising-critical-condition}
    \sinh(2K_x)\sinh(2K_\tau)=1,
\end{equation}
then becomes
\begin{equation}
\label{eq:tfim-critical-from-ising}
    \frac{J}{h}=1,
    \qquad
    \text{or equivalently}\qquad
    h=J.
\end{equation}
In the common normalization $J=1$, this gives the TFIM critical point $h=1$.

\begin{remark}[Why the basis matters]
Many textbooks write the TFIM as
$-J\sum_j\sigma_j^z\sigma_{j+1}^z-\Gamma\sum_j\sigma_j^x$. These notes use the
rotated convention
$-J\sum_j\sigma_j^x\sigma_{j+1}^x-h\sum_j\sigma_j^z$ because it is adapted to the
Jordan--Wigner occupation convention $n_j=(1-\sigma_j^z)/2$. The two forms are
related by a global spin rotation. In the transfer-matrix derivation above, this
rotation is implemented by labeling the classical Ising variable as the eigenvalue
of $\sigma^x$ rather than $\sigma^z$.
\end{remark}

\section{Operator insertions and Euclidean correlators}
\label{sec:operator-insertions-euclidean-correlators}

The transfer-matrix correspondence also maps correlation functions. A diagonal
classical row observable $O(\bm{s})$ becomes an operator $O$ acting on the row
Hilbert space. Its correlation along the transfer direction is
\begin{equation}
\label{eq:transfer-correlator-as-quantum-correlator}
    \langle O_m O_0\rangle_{\rm cl}
    =
    \frac{1}{Z}
    \Tr\left(O T^m O T^{M-m}\right).
\end{equation}
If $T=C\ee^{-a_\tau H}$, then
\begin{equation}
\label{eq:euclidean-quantum-correlator}
    \langle O_m O_0\rangle_{\rm cl}
    =
    \frac{1}{Z_{\rm q}}
    \Tr\left(
        O(\tau)O(0)\ee^{-\beta H}
    \right),
    \qquad
    \tau=m a_\tau,
\end{equation}
where
\begin{equation}
\label{eq:imaginary-time-heisenberg-operator}
    O(\tau)=\ee^{\tau H}O\ee^{-\tau H}.
\end{equation}
At zero temperature, this becomes the ground-state Euclidean correlator. Its
large-$\tau$ decay is controlled by the lowest quantum excitation carrying the
quantum numbers of $O$:
\begin{equation}
\label{eq:euclidean-spectral-decay}
    \langle 0\vert O(\tau)O(0)\vert0\rangle_c
    \sim
    \ee^{-\tau(E_1-E_0)}.
\end{equation}
Thus the transfer-matrix spectrum, the classical correlation length, and the
quantum excitation gap are the same spectral information in different languages.

\section{Quantum-to-classical dictionary}
\label{sec:quantum-classical-dictionary}

For the TFIM/Ising example in the conventions of these notes, the essential
parameter dictionary is
\begin{equation}
\label{eq:tfim-ising-dictionary}
    \boxed{
    K_x=\delta\tau J,
    \qquad
    \ee^{-2K_\tau}=\tanh(\delta\tau h)
    }
\end{equation}
with row transfer matrix
\begin{equation}
\label{eq:tfim-ising-transfer-final}
    T
    \propto
    \exp\left(
        \frac{\delta\tau J}{2}
        \sum_j\sigma_j^x\sigma_{j+1}^x
    \right)
    \exp\left(
        \delta\tau h\sum_j\sigma_j^z
    \right)
    \exp\left(
        \frac{\delta\tau J}{2}
        \sum_j\sigma_j^x\sigma_{j+1}^x
    \right).
\end{equation}
The proportionality constant shifts the quantum free energy by an additive
constant. The operator in the exponential is $-\delta\tau H_{\rm TFIM}$ to
leading order.

The following table summarizes the correspondence.
\begin{center}
\begin{tabularx}{0.95\linewidth}{@{}L{0.31\linewidth}Y Y@{}}
\toprule
Concept & Classical transfer-matrix language & Quantum language \\
\midrule
Transfer direction & Layer index $m$ & Imaginary time $\tau=m a_\tau$ \\
Row configuration & Classical spin row $\bm{s}_m$ & Basis vector of the quantum Hilbert space \\
Transfer matrix & $T_{\bm{s}\bm{s}'}$ & $\ee^{-a_\tau H}$ up to normalization \\
Largest eigenvalue & Dominant Boltzmann weight & Ground-state energy density \\
Eigenvalue ratio & Correlation length $\xi^{-1}=\log(\lambda_0/\abs{\lambda_1})$ & Energy gap $a_\tau(E_1-E_0)$ \\
Classical anisotropy & Different $K_x,K_\tau$ & Spatial velocity and time-lattice spacing \\
Criticality & Divergent correlation length & Closing quantum gap \\
\bottomrule
\end{tabularx}
\end{center}

\section{Hamiltonian limits of commuting transfer matrices}
\label{sec:hamiltonian-limits-commuting-transfer-matrices}

In integrable lattice models one often has a one-parameter family of commuting
transfer matrices,
\begin{equation}
\label{eq:commuting-transfer-matrices}
    [T(u),T(v)]=0,
\end{equation}
where $u$ is a spectral parameter. If $T(u)$ is regular at a point $u=u_0$, the
associated quantum Hamiltonian is commonly obtained from the logarithmic
derivative
\begin{equation}
\label{eq:logarithmic-derivative-transfer-matrix}
    H
    \propto
    \left.
    \frac{\dd}{\dd u}\log T(u)
    \right|_{u=u_0}
    =
    \left.
    T(u)^{-1}\frac{\dd T(u)}{\dd u}
    \right|_{u=u_0},
\end{equation}
plus an additive and multiplicative normalization. Because all $T(u)$ commute,
the Hamiltonian obtained this way commutes with the entire transfer-matrix family:
\begin{equation}
\label{eq:hamiltonian-commutes-transfer-family}
    [H,T(v)]=0.
\end{equation}
This is the transfer-matrix origin of infinitely many conserved quantities in
Bethe-ansatz integrable chains. Later, the six-vertex transfer matrix will yield
the XXZ Hamiltonian in precisely this sense.

\section{Practical summary}
\label{sec:transfer-practical-summary}

The transfer-matrix method should be read as a spectral equivalence principle:
\begin{equation}
\label{eq:transfer-spectral-equivalence-principle}
    \boxed{
    \text{classical layer-to-layer transfer}
    \quad\Longleftrightarrow\quad
    \text{quantum imaginary-time evolution}
    }.
\end{equation}
For the purposes of these notes, the essential workflow is:
\begin{enumerate}[label=(\roman*)]
    \item Write the classical partition function as $Z=\Tr T^M$.
    \item Interpret the layer configuration as a quantum basis state.
    \item If possible, write $T$ as $C\ee^{-a_\tau H}$ or take a controlled
    Hamiltonian limit.
    \item Extract quantum energies from transfer-matrix eigenvalue ratios.
    \item Translate classical correlation lengths into quantum gaps.
    \item Use anisotropic limits to connect two-dimensional classical models to
    one-dimensional quantum Hamiltonians.
\end{enumerate}
The 2D classical Ising model and the 1D TFIM are the canonical example. The same
logic also underlies the relation between the six-vertex model and the XXZ chain,
the eight-vertex model and the XYZ chain, and more general tensor-network
interpretations of Euclidean path integrals.
